\section{Three full arms and an isolated fifth port}
\label{sec:three-arm-isolated-port}

The estimate below does not require positive density for the fourth arm.
It excludes a limiting profile with three full arms and negligible paths from the remaining port.

\begin{theorem}[Three full arms control the four-leaf forest sector]
\label{thm:three-full-arms-forest-control}
There is an absolute constant $C$ with the following property.
Let $T$ be a normalized five-pole of rank $h$, with ports $c,a_1,a_2,a_3,a_4$.
Suppose
\[
 P_{ca_1}(T)=P_{ca_2}(T)=P_{ca_3}(T)=2^h.
\]
Let $F_{1234}(T)$ count forests with odd-degree vertex set exactly
$\{a_1,a_2,a_3,a_4\}$.
Then
\begin{equation}\label{eq:three-arm-forest-control}
 F_{1234}(T)
 \le C\sum_{i=1}^3P_{a_4a_i}(T)+C2^{h/2}.
\end{equation}
In particular, if $h\to\infty$ and all three normalized path counts on the right tend to zero, then $F_{1234}(T)/2^h\to0$.
\end{theorem}

\begin{proof}
Suppress $a_2,a_3,a_4$ and call the resulting two-pole $N$.
Its rank remains $h$, and its $c$--$a_1$ path count remains $2^h$.
No loop can occur, since it would be avoided by every terminal path.
By Lemma~\ref{lem:full-two-pole-classification}, its bridges join saturated bridgeless pieces in series.
Orient this chain from $c$ to $a_1$.
Let $K$ be its final piece, with rank $b$ and right terminal $a_1$.

We first locate the restored ports $a_2,a_3$.
For either port, suppress the other two extra ports and use the local subdivision table in Lemma~\ref{lem:one-full-tripole-sum}.
Suppose its edge belongs to a nonempty capped-ladder piece.
Write $q$ for the suffix position count and $R$ for the total rank of all later pieces.
For an interface or singleton rung, its normalized path count from $c$ equals
\[
 2^{-R}\frac{1+B_q}{2^{q+1}},
\]
where $B_0=1$ and $B_q\le2^q$.
For a square edge, it equals
\[
 2^{-R}\frac{4+3B_q}{2^{q+3}}.
\]
The first expression equals one only when $R=q=0$.
The second is always strictly below one.
For a digon piece, the normalized count is $2^{-R}$.
For a bridge subdivision, it is $2^{-R}<1$, since the right shore has positive rank.

Consequently, both $a_2,a_3$ subdivide edges of $K$.
If $K$ is a nonempty ladder, those edges are right-cap edges or the rung of the final singleton block.
Several restored ports may subdivide the same edge, with their original order retained.
Thus both ports lie within a bounded terminal tail of $K$.

If $b\le5$, the entire piece $K$, with these restored ports, has bounded size.
When $a_4$ is outside this piece, every counted forest contains an $a_4$--$a_i$ path for some $i\le3$.
Its other path lies in $K$ and has only a bounded number of choices.
This proves the first term of~\eqref{eq:three-arm-forest-control}.
When $a_4$ lies inside this piece, both forest paths stay there.
Their count is bounded by an absolute constant.
Increasing $C$ therefore handles every case with $b\le5$.

Assume henceforth that $b\ge6$.
Write the ladder positions of $K$ as $1,\ldots,b-1$.
The left terminal has distance layer zero, and $a_1$ has distance layer $b$.
Each intermediate layer has exactly two vertices.
These are the distance layers from the left terminal: every edge changes its position by at most one, and each rail realizes the stated distance.
Set $d=b-4$.
The tail consisting of layers $d,\ldots,b$ has nine vertices before subdivisions.
It contains both endpoints of every edge containing $a_2$ or $a_3$.

First suppose $a_4$ is outside $K$, or subdivides an edge whose endpoints both lie below layer $d$.
Each counted forest consists of two vertex-disjoint paths.
One joins $a_4$ to a port $a_i$, and the other joins the remaining two leaf ports.
The first path must meet layer $d$.
If the second path visited a layer below $d$, it would use both vertices of layer $d$.
Indeed, its two endpoints lie above that layer, and its two passages are vertex-disjoint.
This is impossible because the first path also needs a vertex there.
Therefore the second path stays in the bounded tail, with the two restored ports included.
The number of its possible edge sets is bounded by an absolute constant.
Summing over $i$ proves
\[
 F_{1234}(T)\le C\sum_{i=1}^3P_{a_4a_i}(T)
\]
in this case.
This reasoning also covers an $a_4$ subdivision of the bridge entering $K$.

It remains to consider an edge containing $a_4$ with at least one endpoint in the tail.
Its other endpoint lies in layer at least $d-1$.
All four leaf ports therefore lie in layers at least $d-1$, with at most three restored vertices.
Every counted forest stays in $K$ with these subdivisions.
A path with both endpoints there cannot leave through a bridge and return.

Here is the required bounded-tail count.
Mark every vertex in layers at least $d-1$, all restored vertices, and the left terminal.
This marks an absolute number of vertices.
If a forest uses an earlier layer, also mark its minimum layer, which adds at most two vertices.
Fix the forest edges incident with marked vertices.
There are only an absolute number of choices for these edge sets.

Choose the path component that reaches the minimum layer $j$.
Both endpoints lie in the marked tail.
In every unmarked layer between $j$ and that tail, this path uses both vertices, with one backward and one forward edge each.
No other path component can use these vertices.
Across each cut between consecutive unmarked layers, the selected edges therefore form a perfect matching.
There are at most two choices.
Two-choice cuts arise only inside square blocks and cannot be consecutive.
For fixed minimum layer $j$, the number of choices is consequently
\[
 O(1)2^{(b-j)/2}.
\]
This bound also covers the constant number of edges beside the marked layers.
Forests wholly inside the marked tail have only $O(1)$ choices.
The geometric sum over all possible minimum layers gives
\[
 F_{1234}(T)
 \le O(1)\left(1+\sum_{j=0}^{b-1}2^{(b-j)/2}\right)
 =O(2^{b/2})\le O(2^{h/2}).
\]
Increasing the same absolute constant $C$ proves the theorem.
\end{proof}

The proof uses the saturated two-pole classification and its exact local subdivision formulas.
It requires no uniform lower bound for any path involving $a_4$.
